% ============================================================
\chapter{Brownian Motion}
\label{ch:brownian}
% ============================================================

In 1827, the botanist Robert Brown observed pollen grains suspended in water through a microscope. They trembled, jiggled, changed direction ceaselessly --- and nobody understood why. It would take until 1905 and a clerk at the Bern patent office, a certain Albert Einstein, for the mystery to begin to clear: these grains were being bombarded by billions of water molecules, and their erratic trajectory reflected the thermal agitation of the fluid. Meanwhile, the mathematician Louis Bachelier had already, in his 1900 thesis, modelled stock market fluctuations using a process with remarkably similar properties. Norbert Wiener, in the 1920s, gave the first rigorous construction. Today, Brownian motion is the cornerstone of stochastic calculus, mathematical finance, and statistical physics --- the continuous-time process from which all others are built.

\section{Definition}

\begin{definition}[Standard Brownian Motion]
\label{def:brownian}
A real-valued stochastic process $(B_t)_{t \geq 0}$ is a \emph{standard Brownian
motion} (or \emph{Wiener process}) if:
\begin{enumerate}
  \item $B_0 = 0$ a.s.,
  \item $(B_t)$ has \emph{independent increments}: for all
    $0 \leq t_1 < t_2 < \cdots < t_n$, the r.v.\
    $B_{t_2} - B_{t_1}, B_{t_3} - B_{t_2}, \ldots, B_{t_n} - B_{t_{n-1}}$
    are independent,
  \item $(B_t)$ has \emph{stationary Gaussian increments}:
    $B_{t+s} - B_s \sim \mathcal{N}(0, t)$ for all $s, t \geq 0$,
  \item the sample paths $t \mapsto B_t(\omega)$ are a.s.\ \emph{continuous}.
\end{enumerate}
\end{definition}

\begin{remark}
Brownian motion is a centred Gaussian process with covariance function
$\Cov(B_s, B_t) = \min(s, t)$.
\end{remark}

\begin{definition}[Multidimensional Brownian Motion]
A $d$-dimensional Brownian motion is $B_t = (B_t^{(1)}, \ldots, B_t^{(d)})$
where $B^{(1)}, \ldots, B^{(d)}$ are independent standard Brownian motions.
\end{definition}

\section{Existence: L\'evy--Ciesielski Construction}

\begin{theorem}[Existence of Brownian Motion]
\label{thm:brownian_existence}
There exists a stochastic process satisfying the four properties of
Definition~\ref{def:brownian}.
\end{theorem}

\begin{proof}[L\'evy--Ciesielski construction (sketch)]
We use the Schauder basis of $C([0,1])$. Let $(Z_n)_{n \geq 0}$ be i.i.d.\
$\mathcal{N}(0,1)$. Define:
\[
  B_t = \sum_{n=0}^{\infty} Z_n \, s_n(t),
\]
where the $s_n$ are the Schauder functions (dyadic hat functions). The series
converges uniformly a.s.\ in $C([0,1])$ and the limit is a Brownian motion
on $[0,1]$. Extension to $\R_+$ is by concatenation.
\end{proof}

\section{Fundamental Properties}

\begin{proposition}[Martingale Properties]
Brownian motion satisfies:
\begin{enumerate}
  \item $(B_t)$ is a martingale with respect to its natural filtration.
  \item $(B_t^2 - t)$ is a martingale.
  \item $\exp(\theta B_t - \theta^2 t/2)$ is a martingale for all $\theta \in \R$.
\end{enumerate}
\end{proposition}

\begin{proposition}[Invariance Properties]
\label{prop:invariance}
The following processes are also standard Brownian motions:
\begin{enumerate}
  \item \textbf{Symmetry}: $(-B_t)_{t \geq 0}$.
  \item \textbf{Scaling}: $(c^{-1/2} B_{ct})_{t \geq 0}$ for any $c > 0$.
  \item \textbf{Time translation}: $(B_{t+s} - B_s)_{t \geq 0}$ for any $s \geq 0$.
  \item \textbf{Time inversion}: $(tB_{1/t})_{t > 0}$ (with $B_{1/t}\big|_{t=0} = 0$).
\end{enumerate}
\end{proposition}

\begin{proposition}[Strong Markov Property]
For every a.s.\ finite stopping time $\tau$, the process
$(B_{\tau + t} - B_{\tau})_{t \geq 0}$ is a Brownian motion independent of
$\mathcal{F}_{\tau}$.
\end{proposition}

\section{Path Regularity}

\begin{theorem}[H\"older Continuity]
\label{thm:holder}
The sample paths of Brownian motion are a.s.\ H\"older continuous with exponent
$\gamma$ for every $\gamma < 1/2$, but not for $\gamma = 1/2$.
\end{theorem}

\begin{theorem}[Nowhere Differentiability]
\label{thm:nowhere_diff}
The sample paths of Brownian motion are a.s.\ \emph{nowhere differentiable}.
More precisely, for almost every $\omega$, there is no $t \geq 0$ such that
$\lim_{h \to 0} (B_{t+h}(\omega) - B_t(\omega))/h$ exists (even as an infinite
limit).
\end{theorem}

\begin{theorem}[Quadratic Variation]
\label{thm:quad_var_bm}
For any sequence of partitions $\Pi_n = \{0 = t_0^n < t_1^n < \cdots < t_{k_n}^n = t\}$
of $[0, t]$ with $\abs{\Pi_n} \to 0$:
\[
  \sum_{i=0}^{k_n - 1} (B_{t_{i+1}^n} - B_{t_i^n})^2 \xrightarrow{L^2} t.
\]
That is, $\langle B \rangle_t = t$.
\end{theorem}

\begin{proof}
Set $Q_n = \sum_{i=0}^{k_n-1} (B_{t_{i+1}} - B_{t_i})^2$. Then
$\E[Q_n] = \sum_i (t_{i+1} - t_i) = t$. Moreover:
\begin{align*}
  \Var(Q_n) &= \sum_i \Var\bigl((B_{t_{i+1}} - B_{t_i})^2\bigr)
  = \sum_i 2(t_{i+1} - t_i)^2 \\
  &\leq 2 \abs{\Pi_n} \sum_i (t_{i+1} - t_i) = 2\abs{\Pi_n} \cdot t \to 0.
\end{align*}
Hence $Q_n \to t$ in $L^2$.
\end{proof}

\begin{theorem}[Infinite Total Variation]
\label{thm:total_var}
The sample paths of Brownian motion a.s.\ have \emph{infinite total variation}
on every interval $[0, t]$ with $t > 0$:
\[
  \sup_{\Pi} \sum_i \abs{B_{t_{i+1}} - B_{t_i}} = +\infty \quad \text{a.s.}
\]
\end{theorem}

\begin{warning}[title=\textbf{Finite Quadratic Variation, Infinite Total Variation}]
Brownian motion has finite quadratic variation ($= t$) but infinite total
variation. This property prevents defining the stochastic integral
$\int_0^t f \, dB$ in the Stieltjes sense and necessitates It\^o's special
construction.
\end{warning}

\section{Law of the Iterated Logarithm}

\begin{theorem}[Law of the Iterated Logarithm]
\label{thm:LIL}
\[
  \limsup_{t \to \infty} \frac{B_t}{\sqrt{2t \log \log t}} = 1 \quad \text{a.s.}
\]
\[
  \liminf_{t \to \infty} \frac{B_t}{\sqrt{2t \log \log t}} = -1 \quad \text{a.s.}
\]
There is an analogous result as $t \to 0$:
\[
  \limsup_{t \to 0^+} \frac{B_t}{\sqrt{2t \log \log(1/t)}} = 1 \quad \text{a.s.}
\]
\end{theorem}

\section{Reflection Principle and Law of the Maximum}

\begin{theorem}[Reflection Principle]
\label{thm:reflection}
Let $\tau_a = \inf\{t \geq 0 : B_t = a\}$ for $a > 0$. For all $b \leq a$:
\[
  \PP(\tau_a \leq t,\; B_t \leq b) = \PP(B_t \geq 2a - b).
\]
\end{theorem}

\begin{corollary}[Law of the Maximum]
\label{cor:law_max}
Let $M_t = \max_{0 \leq s \leq t} B_s$. Then:
\[
  \PP(M_t \geq a) = 2\PP(B_t \geq a) = 2\bigl(1 - \Phi(a/\sqrt{t})\bigr),
  \quad a \geq 0,
\]
where $\Phi$ is the c.d.f.\ of $\mathcal{N}(0,1)$. Moreover,
$M_t \overset{\mathcal{L}}{=} \abs{B_t}$.
\end{corollary}

\begin{corollary}[Distribution of $\tau_a$]
For $a > 0$:
\[
  \PP(\tau_a \leq t) = 2\bigl(1 - \Phi(a/\sqrt{t})\bigr), \quad
  f_{\tau_a}(t) = \frac{a}{\sqrt{2\pi t^3}} \exp\Bigl(-\frac{a^2}{2t}\Bigr),
  \quad t > 0.
\]
In particular, $\E[\tau_a] = +\infty$.
\end{corollary}

\section{Brownian Motion and PDEs}

\begin{theorem}[Feynman--Kac Formula (Simple Case)]
\label{thm:feynman_kac}
Let $g : \R \to \R$ be bounded and continuous. The function
\[
  u(t, x) = \E_x[g(B_T)] = \E[g(x + B_{T-t})], \quad 0 \leq t \leq T,
\]
is the unique bounded solution of the heat equation:
\[
  \frac{\partial u}{\partial t} + \frac{1}{2} \frac{\partial^2 u}{\partial x^2} = 0,
  \quad u(T, x) = g(x).
\]
\end{theorem}

\section{Local Times}

\begin{definition}[Local Time]
The \emph{local time} of Brownian motion at level $a \in \R$ is the process
$(L_t^a)_{t \geq 0}$ defined by the \emph{occupation density formula}:
\[
  \int_0^t f(B_s) \, ds = \int_{-\infty}^{+\infty} f(a) L_t^a \, da
\]
for every Borel function $f \geq 0$. Equivalently:
\[
  L_t^a = \lim_{\varepsilon \to 0} \frac{1}{2\varepsilon}
  \int_0^t \mathbf{1}_{\{\abs{B_s - a} < \varepsilon\}} \, ds.
\]
\end{definition}

\begin{theorem}[Tanaka's Formula]
\label{thm:tanaka}
For all $a \in \R$:
\[
  \abs{B_t - a} = \abs{a} + \int_0^t \text{sgn}(B_s - a) \, dB_s + L_t^a,
\]
where $\text{sgn}(x) = \mathbf{1}_{\{x > 0\}} - \mathbf{1}_{\{x \leq 0\}}$.
\end{theorem}

\section{Exercises}

\begin{exercise}
Show that Brownian motion is a.s.\ non-monotone on every interval.
\end{exercise}

\begin{exercise}
Compute $\E[B_s B_t]$ for $s, t \geq 0$ directly from the definition.
\end{exercise}

\begin{exercise}
Verify that $(t B_{1/t})_{t > 0}$ is a Brownian motion (with the convention
$= 0$ at $t = 0$).
\end{exercise}

\begin{exercise}
Use the reflection principle to compute $\PP(\tau_1 \leq 4)$.
\end{exercise}

\begin{exercise}
Let $(B_t)$ be a Brownian motion and $f \in C^2(\R)$ with $f'' \leq 0$. Show
that $(f(B_t))$ is a local supermartingale.
\end{exercise}

\begin{exercise}
Prove that $\E[\tau_a] = +\infty$ for all $a \neq 0$ using the density of $\tau_a$.
\end{exercise}

\begin{exercise}
Show that the Brownian bridge $X_t = B_t - tB_1$ ($0 \leq t \leq 1$) is a
Gaussian process and compute its covariance function. Verify that
$X_0 = X_1 = 0$ a.s.
\end{exercise}
